Repaso

Excusas

じゅうたい + おくれます = じゅうたいで、おくれます
みちにまよいます + おくれます =　道にまよって、おくれます

La forma te cumple para dar una excusa.

\begin{enumerate}
  \item ニュースを聞いてびっくりしました。
  \item 昨日、夜１１時まで勉強してとてもつかれました。
  \item 待ち合わせの時間におくれてすみません。
  \item テステがじんぱおで
\end{enumerate}

Ir a hacer una accion

Noun + に行きます/来ます
V(quitar ます) + に行きます/来ます

買い物に行きます。
見に行きます

\begin{enumerate}
  \item 図書館に本を借りに行きます。
  \item 日本に１年いて、いろいろな場所にりょこうしに行きたいです。
  \item ゆうびんきょくにきってを買いに行きます。
  \item デパートにかいものに行きます。
  \item あついですから、プールにおよぎに行きます。
\end{enumerate}

Decir que algo pasa antes o despues de algo

Noun + の前に、。。。
Noun + の後で、。。。

Solo Nouns, sustantivos

うんどうの前に朝食をたべます。
うんどうの後で、。。。

Pedir algo con くださいませんか

V(en forma て)くださいませんか

来週の木曜日に来てくださいませんか。


Preguntar si ya se hizo una accion

もう＿V-ました(か)。

もう朝ごはんを食べましたか。
もうアベンジャースのえいがを見ませんか。

Describir propiedad de un tema
N1はN2が イA/ナA です。

中国語ははつおんはむずかしいです。
ミゲル先生は背が高いです。

% すめたい Frio personalidad

\begin{enumerate}
  \item フタバの学校は宿題がかんたんです。
  \item 日本語の先生ははつおんがとくいです。
  \item 日本はりょうりがうまいです。
  \item メキシコは天気がたいへんです。
  \item ドイツ人は背が高いです。
  \item メキシコ人はせいかくが面白いです。
\end{enumerate}

Preguntar por una cosa y decir si esta bien o no (sugerir alternativa)

Sustantivoでもいいですか。
ええ、いいですよ。
ええ、大丈夫です。

今日、かいぎはちょっと。。。
あしたでもいいですか。

Describir propiedad de un tema

NはVるのがイA/ナAです。

Decir la forma de hacer algo con かた

V(quitar ます)+かた

食べ方 La forma de comer
作り方 La forma de hacer
書き方 La forma de escribir

海
